% !TeX root = ../main.tex

\subsection{Threats and Limitations}

\textbf{Threats.}\todel{The threats to validity are subjected to three aspects. } First,~as the campaign against malware continues, new strains of malicious packages emerge daily. This implies that the original dataset of malicious packages used in our evaluation may not encompass the full spectrum of new malicious packages, hindering \tool's effectiveness. However, we have demonstrated the effectiveness of \tool in learning from newly published packages and adapting to detect emerging threats. Second, while we thoroughly examine~the~potentially malicious package versions identified by \tool, we acknowledge that~the~vast~workload makes it impractical to validate every potentially benign package version. Therefore, we only systematically report the result of precision but not the result of recall. Third, \tool can generate false positives. To address this, we conduct manual validation on identified package versions before reporting them to PyPI and NPM teams. While this process can also introduce false positives, the fact that the reported packages are subsequently removed indicates the validity of our validation. 
\tosemrevision{Fourth, as pre-trained language models continue to advance rapidly, we expect to see new and more powerful models. Nevertheless, the choice of language models is orthogonal to our approach. We plan to test the effectiveness of additional language models on this task.}
\tosemrevision{Fifth, the varying behavior sequences between malicious and benign packages may reduce the accuracy of \tool in cross-lingual scenarios. Nevertheless, the cross-lingual setting represents an extreme case in real-world where there are no malicious samples within the same ecosystem.}


% This limitation could potentially result in false negatives, where some malicious packages are mistakenly classified as benign. 


\new{\textbf{Limitations.} \tosem{First, our method is designed for PyPI and NPM packages. However, our method can be easily generalized into other ecosystems by implementing the feature extractor and behavior sequence generator for a new ecosystem. The feature set is language-agnostic. By leveraging the existing AST parser and call graph generator of the new ecosystem, we believe that the integration can be straightforward.} Second, large packages may have a large size of the behavior sequence which can exceed the 512 token limit of BERT. However, based on our analysis of packages in Table \ref{table:monitor-stats}, we identify a relatively modest proportion (\todo{13.6\%}) of packages that exceed the limit. We plan to split the behavior sequence into multiple segments to circumvent the limit while maintaining their sequential integrity.}
\tosemrevision{Third, the proposed approach encounters higher false positives in real-world setting. We foresee several challenges and opportunities that could help reduce false positive rates. On the one hand, we can leverage ensemble methods. Utilizing a combination of different detection tools could be advantageous. An ensemble approach, where decisions are based on the consensus among multiple tools, might improve accuracy and reduce false positive rates. On the other hand, we can utilize dynamic analysis during our detection. Incorporating dynamic analysis into the detection process could provide deeper insights into the behavior of packages by capturing detailed runtime behaviors, potentially reducing false positives by confirming suspicious activities through actual execution rather than static analsysis. Nonetheless, with a large number of new packages being released daily, it remains an open question how to leverage the benefits of dynamic analysis while mitigating its computational and time costs.}

% The limitations of \tool is two-fold.
% First, the features are not exhaustive.

% Second, the unsound call graph generation can hinder the building of subsequent behavior sequences, causing \tool fails to interpret malicious knowledge from features.




% \section{Discussion}

% We have 
% \textbf{Future work.} interpreted languages. 

% \textbf{}

% Implication

% 语言扩展
%     解释性
%     非解释型


% 定位 
